We've climbed parts of the ladder through the experiments. This chapter
pulls together what we found into the picture the ladder was always
meant to give us: a map of where capability appears, where it holds, and
where it breaks.

The ladder is a staircase of increasingly demanding skills. We can see
how many steps the model climbed --- and, more interestingly, the exact
step where it stopped, and what it took to get there.

The funny thing about building a ladder is that the interesting
discovery is rarely how high we can climb. It is finding the precise
rung where the footing gives out, and trying everything we can to get
past it --- so that we know, with real confidence, where the boundary
lies.

\section{The ladder, with the measured
results}\label{the-ladder-with-the-measured-results}

\begin{longtable}[]{@{}lll@{}}
\toprule\noalign{}
Level & Capability & Measured result \\
\midrule\noalign{}
\endhead
\bottomrule\noalign{}
\endlastfoot
1 & arithmetic & ✅ learned, \textasciitilde0.90 in-dist \\
2--3 & variables, conditionals & ✅ learned \\
4--5 & lists, indexing & ✅ learned \\
6 & loops & ✅ learned \\
7 & \textbf{composition} & ❌ \textbf{0.000} --- the wall \\
8 & simple functions & not isolated \\
\end{longtable}

The model climbed the early rungs cleanly. It stopped at
\textbf{composition}.

\section{Where capability appears, holds, and
breaks}\label{where-capability-appears-holds-and-breaks}

\textbf{Appears:} individual operations are learned reliably
(\textasciitilde0.90 in-distribution) and generalize across
\emph{instances} (held-out + canary → general program). The model is a
genuine learner for a single operation.

\textbf{Holds only within a surface form.} The moment the surface
changes --- different variable names, or a different way of saying the
same instruction --- the capability collapses (0.350 on new identifiers,
\textbf{0.000} on new wording). So the early rungs are not
\emph{robust}; they are \emph{phrasing-tied}.

\textbf{Breaks at composition.} The most important boundary: the model
cannot combine two things it knows. \texttt{abs} + \texttt{add} → 0.000,
even though it knows both. And this break is \textbf{invariant to scale}
--- it survives across the full size range we tested.

\section{The model-size and data-size
boundaries}\label{the-model-size-and-data-size-boundaries}

This is where the breadth of the brew matters. We did not test just one
bigger model; we tested four, and we also scaled the data. Here are the
real numbers.

\begin{itemize}
\tightlist
\item
  \textbf{Model size:} we ran the composition test across a
  \textasciitilde33× range of capacity --- a 0.88M pilot, a 6.5M model,
  a 10.2M model, and a 29.3M model. All four scored \textbf{0.000} on
  composition. (Their in-distribution scores were 0.900, 0.900, 0.383,
  and 0.897 respectively --- a side note we return to honestly, because
  the 10.2M model actually dropped in-distribution, showing that bigger
  is \emph{not} monotonically better at this scale.) The composition
  wall is not a capacity ceiling in this regime.
\item
  \textbf{Data size:} 100 → 3000 examples improved in-distribution (0.10
  → 0.64) but did not touch composition. More quantity helps the shallow
  score, not the deep boundary.
\end{itemize}

\section{The failure map}\label{the-failure-map}

\begin{longtable}[]{@{}
  >{\raggedright\arraybackslash}p{0.3325\linewidth}
>{\raggedright\arraybackslash}p{0.6175\linewidth}
@{}}
\toprule\noalign{}
\begin{minipage}[b]{\linewidth}\raggedright
Failure mode
\end{minipage} & \begin{minipage}[b]{\linewidth}\raggedright
On what the model breaks
\end{minipage} \\
\midrule\noalign{}
\endhead
\bottomrule\noalign{}
\endlastfoot
Surface fragility & new wording → 0.000; new identifiers → 0.350 \\
Composition & unseen combination of known primitives → 0.000 \\
Continual forgetting & sequential A→B destroys A (0.000); replay fixes
(1.000) \\
\end{longtable}

\section{What the ladder finally tells
us}\label{what-the-ladder-finally-tells-us}

The ladder is not a shelf of ``caps'' the model unlocks. It is a map of
\emph{where the boundary lies}. For our \textasciitilde0.9M model:

\begin{itemize}
\tightlist
\item
  it genuinely climbs the individual-capability rungs,
\item
  but each rung is \textbf{surface-tied}, not robust,
\item
  and the ladder stops at \textbf{composition} --- the one thing that
  separates ``a bag of reflexes'' from ``capability.''
\end{itemize}

And the most useful part: this map is honest about \emph{what doesn't
change it}. The composition boundary survived a \textasciitilde33×
capacity increase and a 30× data increase. That is a strong, falsifiable
statement about where the limit lives and where it does \emph{not} live
--- which points the next investigation at \emph{structural} data and
inductive bias, not at scaling.

There is also a quiet lesson in the shape of the map. The model is not
uniformly good or uniformly bad. It is good at one thing --- reproducing
a single operation it was shown --- and bad at everything that requires
putting operations together. That shape is far more informative than a
single performance number, because it tells us exactly where nature of
the limitation lives and where it does not.
